\begin{prop}
    \label{prop:stubCharacterization}
    Let $(W,S)$ be an arbitrary Coxeter system and let $w$ be an FC element in $W$. Let $w=w_p\dots w_2w_1$ be the
    Cartier--Foata form of $w$. Then $w$ is a left stub if and only if every generator $t\in \supp(w_2)$ fails to
    commute with at least two generators $s,s'\in \supp(w_1)$. 
\end{prop}

\begin{remark}
    Note that the above condition on $t$ is vacuously true when $w=w_1$, i.e. when $w$ is a product of commuting
    generators.
\end{remark}

\begin{proof}[Proof of Proposition~\ref{prop:stubCharacterization}]
    Let $t\in \supp(w_2)$ and $s\in \supp(w_1)$. Then $s\neq t$ by
    Proposition~\ref{prop:FCCriterion}.(1), so $t$ is covered by $s$ in $H(w)$ if and only if $m(s,t)\ge 3$.  The
    proposition then follows from Remark~\ref{rmk:heapsAndWords} and Proposition~\ref{prop:Reducibility}.
\end{proof}
\noindent 
The proposition shows that whether an FC element is a left stub is determined ``locally'', by only the first two
layers of its Cartier--Foata form. 

\begin{corollary}
    \label{coro:firstTwoLayers}
    Let $(W,S)$ be an arbitrary Coxeter system.  Let $G$ be the Coxeter diagram of $(W,S)$.  Then an
    $\la(2)$-element $w\in W$ is a left stub if and only if in the Cartier--Foata form $w=w_p\dots w_2w_1$ of $w$,
    we have 
    \begin{enumerate} 
        \item  $w_1=ss'$ for two commuting generators $s, s'$ in $S$;
        \item Every generator $t\in \supp(w_2)$ is adjacent to both $s$ and $s'$ in $G$.  
    \end{enumerate}
\end{corollary}

\begin{proof}
    By Proposition~\ref{prop:stubCharacterization}, it suffices to show that $\supp(w_1)$ contains exactly two
    generators.  By Corollary~\ref{coro:layerSizeBound}, it further suffices to show that
    $\supp(w_1)$ is not empty or a singleton.  This follows from Proposition~\ref{prop:facts01}: if $w_1$ is
    empty, then $w=1$ and $\la(w)=\la(1_W)=0$; if $\supp(w_1)$ contains only one generator $s$, then
    $w_2$ must be empty by Proposition~\ref{prop:stubCharacterization}, which forces $w=w_1=s$ and
    $\la(w)=\la(s)=1$.    
\end{proof}

For nontrivially $\la(2)$-finite systems, Corollary~\ref{coro:firstTwoLayers} turns out to impose strong
restrictions on the first two layers of $\la(2)$-stubs.  The following observation imposes further
restrictions on the deeper layers. 

\begin{lemma}
    \label{lemm:deeperLayers}
    Let $(W,S)$ be an arbitrary Coxeter system. Let $G$ be the Coxeter diagram
    of $(W,S)$. Let $w\in \fc(W)$, let $w=w_p\dots w_1$ be the Cartier--Foata form of $w$, and let~$2<i\le p$. 
    \begin{enumerate}
        \item Every element $t\in \supp(w_i)$ is adjacent in $G$ to some element $s\in \supp(w_{i-1})$.  
        \item If $s,t\in S$ are generators connected by a simple edge in $G$ such that $s$ appears in $w_{i-2}$ and
             $w_{i-1}=t$, then $s$ does not appear in $w_i$.  
    \end{enumerate}
\end{lemma}

\begin{proof}
    Part~(1) simply paraphrases Condition (2) from the characterization of the Cartier--Foata form
    in \autoref{sec:FC}. To prove (2), note that if $s$ appears in $w_{i}$ then we may commute the
    $s$ in $w_{i-2}$ to the left and the $s$ in $w_i$ to the right, if necessary, until they are
    both next to the $t$ in $w_{i-1}$. This results in a reduced factorization $w=x \cdot sts\cdot
    y$ of $w$, contradicting the word criterion for full commutativity. 
\end{proof}

The following example shows in detail how Conditions (1)--(2) from Corollary~\ref{coro:firstTwoLayers} and
Lemma~\ref{lemm:deeperLayers} strongly restrict the Cartier--Foata forms of left $\la(2)$-stubs in a
nontrivially $\la(2)$-finite Coxeter system.  Roughly speaking, the restrictions are strong because
the Coxeter diagram of such a system always contains few pairs of vertices that share a neighbor as
well as few heavy edges.

\begin{example}
    \label{eg:layerAnalysis} 
    We find all left $\la(2)$-stubs in the Coxeter system $(W,S)$ of type $B_n$ in this example.
    Table~\ref{tab:B4stubs} contains the heaps of the six stubs in $\cals(B_4)$, which will serve as an example. 

    Let $w\in\sw$ and let $w=w_p\dots w_2w_1$ be the Cartier--Foata form of $w$.  By Condition (1)
    of Corollary~\ref{coro:firstTwoLayers}, we must have either $w_1=ij$ for some $1\le i,j\le n$ where
    $j>i+2$ or $w_1=i(i+2)$ for some $1\le i\in n-1$. In the former case---which holds for only the
    stub in the third column in Table~\ref{tab:B4stubs} if $n=4$---the second layer of $w$ cannot exist
    by Condition (2) of Corollary~\ref{coro:firstTwoLayers} since $i,j$ share no neighbor in $G$.  In the
    latter case---which corresponds to the first two columns of Table~\ref{tab:B4stubs}---either $w_2$
    does not exist and $w=w_1=i(i+2)$ (first row), or $w_2$ exists and is given by $w_2=(i+1)$, the
    only common neighbor of $i$ and $(i+2)$ in $G$ (second row). In the latter subcase, if $i>1$
    then we have $m(i+1,i)=m(i+1,i+2)=3$, hence Lemma~\ref{lemm:deeperLayers} rules out the possibility
    of a third layer in $w$ and forces $w=w_2w_1=(i+1)\cdot i(i+2)$; if $i=1$, then
    Lemma~\ref{lemm:deeperLayers} implies that we have either $w=w_{2}w_1=2\cdot (13)$ or
    $w=w_3w_2w_1=1\cdot 2\cdot (13)$. In particular, if $w$ contains at least three layers then we
    have to have $w_1=13, w_2=2$ and $w_3=1$, whence $w$ cannot have a fourth layer $w_4$: an
    element $t\in \supp(w_2)$ has to be adjacent to the only generator $1$ in $w_3$ and thus has to
    be $t=2$, so if $w$ contains a fourth layer then $w=\dots\cdot 2\cdot 1\cdot 2\cdot
    (13)=\dots\cdot (2121)\cdot 3$, contradicting the fact that $w$ is FC. The stub $1\cdot 2\cdot
    (13)$ is shown in the third row and first column in Table~\ref{tab:B4stubs}.

    In summary, we can find all left $\la(2)$-stubs in $B_n$, and among them there is a unique stub
    with more than two layers, namely $1\cdot 2\cdot (13)$. The problem of finding left
    $\la(2)$-stubs in the Coxeter system $A_n$ is similar but easier.
\end{example}

\begin{table}
    \begin{tabular}{|c|c|c|}
        \hline
        \begin{tikzpicture}
            \node[main node] (1) {};
            \node[main node] (2) [right=1cm of 1] {};
            \node (11) [above = 0.1cm of 1] {\tiny{$1$}};
            \node (22) [above = 0.1cm of 2] {\tiny{$3$}};
        \end{tikzpicture}
        &
        \begin{tikzpicture}
            \node[main node] (1) {};
            \node[main node] (2) [right=1cm of 1] {};
            \node (11) [above = 0.1cm of 1] {\tiny{$2$}};
            \node (22) [above = 0.1cm of 2] {\tiny{$4$}};
        \end{tikzpicture} 
        &
        \begin{tikzpicture}
            \node[main node] (1) {};
            \node[main node] (2) [right=1cm of 1] {};
            \node (11) [above = 0.1cm of 1] {\tiny{$1$}};
            \node (22) [above = 0.1cm of 2] {\tiny{$4$}};
        \end{tikzpicture}
        \\
        \hline
        \begin{tikzpicture}
            \node[main node] (1) {};
            \node[main node] (2) [right=1cm of 1] {};
            \node[main node] (3) [below right = 0.5cm and 0.5cm of 1] {};
            \node (11) [above = 0.1cm of 1] {\tiny{$1$}};
            \node (22) [above = 0.1cm of 2] {\tiny{$3$}};
            \node (33) [below = 0.1cm of 3] {\tiny{$2$}};
            \path[draw]
            (1)--(3)--(2); 
        \end{tikzpicture}
        &
        \begin{tikzpicture}
            \node[main node] (1) {};
            \node[main node] (2) [right=1cm of 1] {};
            \node[main node] (3) [below right = 0.5cm and 0.5cm of 1] {};
            \node (11) [above = 0.1cm of 1] {\tiny{$2$}};
            \node (22) [above = 0.1cm of 2] {\tiny{$4$}};
            \node (33) [below = 0.1cm of 3] {\tiny{$3$}};
            \path[draw]
            (1)--(3)--(2);  
        \end{tikzpicture}
        &        
        \multicolumn{1}{c}{}\\
        \cline{1-2}
        \begin{tikzpicture} 
            \node[main node] (1) {};
            \node[main node] (2) [right=1cm of 1] {};
            \node[main node] (3) [below right = 0.5cm and 0.5cm of 1] {};
            \node[main node] (4) [below = 1cm of 1] {};
            \node (11) [above = 0.1cm of 1] {\tiny{$1$}};
            \node (22) [above = 0.1cm of 2] {\tiny{$3$}};
            \node (33) [below = 0.1cm of 3] {\tiny{$2$}};
            \node (11) [below = 0.1cm of 4] {\tiny{$1$}};
            \path[draw]
            (1)--(3)--(2)
            (3)--(4);
        \end{tikzpicture}            
        &
        \multicolumn{2}{c}{}\\
        \cline{1-1}
    \end{tabular}
    \vspace{1em} 
    \caption{The six stubs of $B_4$}
    \label{tab:B4stubs}
\end{table}

We classify the left $\la(2)$-stubs of all nontrivially $\la(2)$-finite Coxeter systems in the
following theorem.  In the theorem, we express every element in its Cartier--Foata form and use
$\cdot$ to separate the different layers, and the symbol $\sqcup$ denotes disjoint union. We will
draw the heaps of typical stubs immediately after the proof of the theorem.  The classification lays
the foundation for the subsequent study of $\la(2)$-cells. 

\begin{thm}
    \label{thm:stubDescription}
    Let $(W,S)$ be an irreducible and nontrivially $\la(2)$-finite Coxeter system. Suppose that $(W,S)$ is of type
    $X$, and let $G$ be the Coxeter diagram of $(W,S)$, labeled in the convention explained in Remark~\ref{rmk:labels}.
    Let $\sx$ be the set of all left stubs of $\la$-value 2 in $W$.  Then $\sx$ can be described as follows. 
    \begin{enumerate}
        \item If $X=A_n$ for some $n\ge 3$, then $\sx=\mathcal{S}_1\sqcup\mathcal{S}_2$ where 
            \[
            \mathcal{S}_1=\{x_{ij}:=ij\,\vert\, 1\le i,j\le n,j>i+1\}\]
            and
            \[
                \mathcal{S}_2=\{y_i:=i\cdot
                (i-1)(i+1)\,\vert\, 1< i<n\}.
            \]
        \item If $X=B_n$ for some $n\ge 3$, then $\sx=\mathcal{S}_1\sqcup\mathcal{S}_2\sqcup \cals_3$ where
            $\mathcal{S}_1,\mathcal{S}_2$ are the same \textup{(}see Remark
           ~\ref{rmk:stubRemarks1}.\textup{(}1\textup{)}\textup{)} as in \textup{(}1\textup{)}
            and 
            \[
                \cals_3=\{z_1:=1\cdot 2\cdot 13\}.
            \] 
        \item If $X=\tilde C_{n-1}$ for some $n\ge 5$, then $\sx =\cals_1\sqcup \cals_2\sqcup
             \cals_3\sqcup\cals_4$ where $\cals_1,\cals_2,\cals_3$ are the same as in
             \textup{(}2\textup{)} and 
            \[
                \cals_4=\{z_n:=n\cdot (n-1)\cdot (n-2)n\}.
            \]
        \item If $X=E_{q,r}$ for some $r\ge q\ge 1$, then $\sx=\cals_1\sqcup\cals_2\sqcup
             \mathcal{T}_2\sqcup\cals_3$ for the sets 
            \[
            \cals_1=\{x_{st}\,\vert\,s,t\in S, m(s,t)=2\},\]
            \[
                \cals_2=\{y_{i}:= i\cdot (i-1)(i+1)\,\vert\, -q< i<r\},
            \]
            \[
                \mathcal{T}_2=\{y'_0:=0\cdot (-1)v\}\sqcup\{y''_0:=0\cdot 1v\},
            \]
            and $\cals_3=\{z_s: s\in S, s\neq 0\}$ where 
            \[
                z_s= 
                \begin{cases}
                    s\cdot (s-1)\cdot \ldots \cdot 0\cdot (-1)v &\text{if $s\neq v$ and $1\le s\le
                    r$},\\
                    s\cdot (s+1)\cdot \ldots \cdot 0\cdot (1)v & \text{if $s\neq v$ and $-q\le s\le -1$},\\
                    v\cdot 0 \cdot (-1)1&\text{if $s=v$}.
                \end{cases}
            \]
        \item If $X=F_n$ for some $n\ge 4$, then $\sx=\cals_1\sqcup\cals_2\sqcup\cals_3$
            where  $\cals_1,\cals_2$ are the same as in \textup{(}1\textup{)} and 
            \[
                \cals_3=\{
                    z_1:=1\cdot 2\cdot 3\cdot 24, z_2:=2\cdot 3\cdot
                24\}\sqcup\{z_i\,\vert\, 3\le i\le n\}, 
            \]
            where $z_i=i\cdot (i-1)\cdot \ldots\cdot 3\cdot 2\cdot 13$ for each $3\le i\le n$.
        \item If $X=H_n$ for some $n\ge 3$, then $\sx=\cals_1\sqcup\cals_2\sqcup\cals_3\sqcup\cals_4$ where
            $\cals_1,\cals_2,\cals_3$ are the same as in \textup{(}2\textup{)} and 
            \[
                \cals_4=\{z_i:=i\cdot (i-1)\cdot \ldots\cdot 1\cdot 2\cdot 13\,\vert\,
                2\le i\le n\}.
            \]
        \item In each of the above, we listed each stub of $\la$-value 2 exactly once. {\textup{(}}For example, two
            $x_{ij},x_{i'j'}$ from $ \cals_1$ are equal only if $i=i',j=j'$.{\textup{)}} The cardinality of $\sx$ 
            in $W$ is given by 
            \[ 
                \abs{\sx}=
                \begin{cases}\binom{n}{2} -1 & \text{if $X=A_n$ \textup{(}$n\ge 3$\textup{)}};\\ 
                    \binom{n}{2} & \text{if
                      $X=B_n$ \textup{(}$n\ge 3$\textup{)}};\\ 
                      \binom{n}{2} +1 & \text{if $X=\tilde C_{n-1}$
                      \textup{(}$n\ge 5$\textup{)}};\\ 
                      \binom{n+1}{2}-1 & \text{if $X= E_{q,r}$ \textup{(}$q,r\ge 1, n=q+r+2=\abs{S}$\textup{)}};\\ 
                      \binom{n+1}{2}-1 & \text{if $X= F_n$ \textup{(}$n\ge 4$\textup{)}};\\
                      \binom{n+1}{2}-1 & \text{if $X= H_n$ \textup{(}$n\ge 3$\textup{)}}.
                \end{cases} 
            \]
    \end{enumerate}
\end{thm}


\begin{remark}
    \label{rmk:stubRemarks1}
    \begin{enumerate}
        \item In Theorem~\ref{thm:stubDescription}, by saying that a certain set $\cals_i$ in some part  is the same as
             the set $\cals_i$ given from an earlier part, we mean that the former set consists of the elements
             given by the same reduced words as those of the elements in the latter set. 
        \item By symmetry, an element $w\in W$ is right stub if and only if $w\inverse$ is a left stub, so
             $\cals'(W)=\{w\inverse: w\in \cals(W)\}$ for each Coxeter system $(W,S)$ appearing in Theorem
            ~\ref{thm:stubDescription}.
    \end{enumerate} 
\end{remark}


\begin{proof}[Proof of Theorem~\ref{thm:stubDescription}]
    Part~(7) follows from Parts (1)--(6) by inspection (to see that every element in $\sw$ is indeed
    listed only once) and by easy counting arguments, so it suffices to show that in each of Cases
    (1)--(6), an element $w\in W$ is a left \twostub if and only if $w\in \sx$. The ``only if''
    implication follows from analysis of the layers of left $\la(2)$-stubs:
    Corollary~\ref{coro:firstTwoLayers} implies that the first two layers of such a stub $w$ must satisfy
    Conditions (1)--(2) given in the corollary, whence Lemma~\ref{lemm:deeperLayers} forces $w$ to be an
    element in $\sw$ by arguments similar to those we used for type $B_n$ in
    Example~\ref{eg:layerAnalysis}. The ``if'' implication, i.e. the fact that every element $w\in \sw$ is
    a left $\la(2)$-stub, can be proven as follows: first, draw the heap of the specified reduced
    word of $w$ we gave (such as $x_{13}=13\in A_n$) and use Proposition~\ref{prop:FCCriterion} to see that $w$
    is indeed FC; second, use Proposition~\ref{prop:Reducibility} to see that $w$ is a stub;
    finally, note that we can transform $w$ to its first layer $w_1$ by a sequence of left lower
    star operations by the analog of Proposition~\ref{prop:Reducibility} for left star operations,
    and then invoke Corollary~\ref{coro:starA} to conclude that $\la(w)=\la(w_1)=2$.  The proof is complete.
\end{proof}

Let us describe and draw the heaps of stubs described in Theorem~\ref{thm:stubDescription}. Call a stub
$w\in \sw$ \emph{short} if $l(w)=2$, \emph{medium} if $l(w)=3$, and \emph{long} if $l(w)>3$. Then
the short, medium, and long stubs in $\sw$ are precisely those labeled by ``$x$'', ``$y$''
(including $y'_0$ and $y''_0$ in type $E_{q,r}$) and ``$z$'' in Theorem~\ref{thm:stubDescription}; they are
also precisely those with one, two, and more than two layers, respectively.  The heaps of the short
and medium stubs take the forms shown in Table~\ref{tab:xyStubs}.  Long stubs exist in all
$\la(2)$-finite Coxeter systems except $A_n$.  In type $B,\tilde C, F$ and $H$, every long stub
involves a ``turn'' at the second layer, that is, there is always a heavy edge $\{s,t\}$ in the
Coxeter diagram such that $s$ appears in the first and third layers of the stub while $t$ appears in
the second layer. The stub $1213\in B_4$ pictured in Table~\ref{tab:B4stubs} is such an example.
Table~\ref{tab:zBCFH} contains more typical such stubs.  Finally, in type $E_{q,r}$, each long stub $w$
with Cartier--Foata form $w=w_p\dots w_2w_1$ satisfies the conditions $\supp(w_2)=\{0\}$,
$\supp(w_1)=\{s,t\}$, $\supp(w_2)=\{u\}$ where $\{s,t,u\}=\{-1,1,v\}$; the heaps of the long stubs are
shown in Table~\ref{tab:zE}.

    \begin{table}
        \begin{tabularx}{0.8\textwidth}{|X|X|}
            \hline
            \raisebox{-1.2em}{
    \begin{tikzpicture}
        \node (0) {short stubs:};
        \node [main node] (1) [right = 0.3cm of 0] {};
        \node [main node] (2) [right = 1cm of 1] {};
\end{tikzpicture}
}
    &
    \raisebox{-2em}{
    \begin{tikzpicture}
        \node (00) {medium stubs:};
        \node [main node] (4) [above right = 0.05cm and 0.3cm of 00] {};
        \node [main node] (5) [below right = 0.5cm and 0.5cm of 4] {};
        \node [main node] (6) [above right = 0.5cm and 0.5cm of 5] {};
        \node (x) [below =1.5em of 6] {};
        \path[draw]
        (4)--(5)--(6);
    \end{tikzpicture}
}
    \\
    \hline
\end{tabularx}
\vspace{1em} 
\caption{The short and medium stubs of $\sw$}
    \label{tab:xyStubs}
    \end{table}


    \begin{table}
        \centering
        \begin{tabular}{|c|c|}
            \hline
            \raisebox{0.8em}{
                \begin{tikzpicture}
                    \node (0) {$z_{n}\in \tilde C_{n-1}$:};
                    \node [main node] (1) [right = 0.8cm of 0] {};
                    \node [main node] (2) [above left = 0.5cm and 0.5cm of 1] {};
                    \node [main node] (3) [above right = 0.5cm and 0.5cm of 1] {};
                    \node [main node] (4) [below right = 0.5cm and 0.5cm of 1] {};

                    \node (11) [right = 0.01cm of 1] {\tiny{$(n-1)$}};
                    \node (22) [above = 0.01cm of 2] {\tiny{$(n-2)$}};
                    \node (33) [above = 0.01cm of 3] {\tiny{$n$}};
                    \node (44) [below = 0.01cm of 4] {\tiny{$n$}};
                    \path[draw]
                    (2)--(1)--(3)
                    (1)--(4);
            \end{tikzpicture}}
            & 
            \raisebox{-0.1em}{
                \begin{tikzpicture}
                    \node (x0) {$z_{1}\in F_{n}$:};
                    \node [main node] (x1) [right = 1.3cm of x0] {};
                    \node [main node] (x2) [above left = 0.5cm and 0.5cm of x1] {};
                    \node [main node] (x3) [above right = 0.5cm and 0.5cm of x1] {};
                    \node [main node] (x4) [below left = 0.5cm and 0.5cm of x1] {};
                    \node [main node] (x5) [below left = 0.5cm and 0.5cm of x4] {};

                    \node (x11) [below = 0.01cm of x1] {\tiny{$3$}};
                    \node (x22) [above = 0.01cm of x2] {\tiny{$2$}};
                    \node (x33) [above = 0.01cm of x3] {\tiny{$4$}};
                    \node (x44) [below = 0.01cm of x4] {\tiny{$2$}};
                    \node (x55) [below = 0.01cm of x5] {\tiny{$1$}};
                    \path[draw]
                    (x2)--(x1)--(x3)
                    (x1)--(x4)--(x5);
                \end{tikzpicture} 
            }
            \\ \hline

            \raisebox{0.2em}{
                \begin{tikzpicture}
                    \node (y0) {$z_i\in F_n$:};
                    \node (b) [below = 0.01cm of y0] {\tiny{$(3\le i\le n)$}};
                    \node [main node] (y1) [above right = 0.1cm and 0.8cm of y0] {};
                    \node [main node] (y2) [above left = 0.5cm and 0.5cm of y1] {};
                    \node [main node] (y3) [above right = 0.5cm and 0.5cm of y1] {};
                    \node [main node] (y4) [below right = 0.5cm and 0.5cm of y1] {};
                    \node [main node] (y5) [below right = 0.8cm and 0.8cm of y4] {};
                    \node [main node] (y6) [below right = 0.5cm and 0.5cm of y5] {};

                    \node (y11) [right = 0.01cm of y1] {\tiny{$2$}};
                    \node (y22) [above = 0.01cm of y2] {\tiny{$1$}};
                    \node (y33) [above = 0.01cm of y3] {\tiny{$3$}};
                    \node (y44) [right = 0.01cm of y4] {\tiny{$3$}};
                    \node (y44) [right = 0.01cm of y5] {\tiny{$(i-1)$}};
                    \node (y44) [right = 0.01cm of y6] {\tiny{$i$}};
                    \path[draw]
                    (y2)--(y1)--(y3)
                    (y1)--(y4)
                    (y5)--(y6);
                    \path[draw,dashed]
                    (y4)--(y5);
            \end{tikzpicture}}
            & 
            \begin{tikzpicture}
                \node (yx0) {$z_{i}\in H_{n}$:};
                \node (b) [below = 0.01cm of yx0] {\tiny{$(1\le i\le n)$}};
                \node [main node] (yx1) [right =  0.8cm of yx0] {};
                \node [main node] (yx2) [above left = 0.5cm and 0.5cm of yx1] {};
                \node [main node] (yx3) [above right = 0.5cm and 0.5cm of yx1] {};
                \node [main node] (yx4) [below left = 0.5cm and 0.5cm of yx1] {};
                \node [main node] (yx5) [below right = 0.5cm and 0.5cm of yx4] {};
                \node [main node] (yx6) [below right = 0.8cm and 0.8cm of yx5] {};
                \node [main node] (yx7) [below right = 0.5cm and 0.5cm of yx6] {};

                \node (yx11) [right = 0.01cm of yx1] {\tiny{$2$}};
                \node (yx22) [above = 0.01cm of yx2] {\tiny{$1$}};
                \node (yx33) [above = 0.01cm of yx3] {\tiny{$3$}};
                \node (yx44) [right = 0.01cm of yx4] {\tiny{$1$}};
                \node (yx55) [right = 0.01cm of yx5] {\tiny{$2$}};
                \node (yx66) [right = 0.01cm of yx6] {\tiny{$(i-1)$}};
                \node (yx77) [right = 0.01cm of yx7] {\tiny{$i$}};
                \path[draw]
                (yx2)--(yx1)--(yx3)
                (yx1)--(yx4)--(yx5)
                (yx6)--(yx7);
                \path[draw,dashed]
                (yx5)--(yx6);
            \end{tikzpicture}\\ \hline
        \end{tabular}
        \vspace{1em} 
        \caption{Typical long stubs of $\sw$ in $B_n, \tilde C_{n-1}, F_n$ and $H_n$}
        \label{tab:zBCFH}
    \end{table}

    \begin{table}
        \begin{tabularx}{\textwidth}{|X|c|X|}
            \hline
            \begin{tikzpicture}
                \node (y0) {$z_i$:};
                \node (b) [below = 0.01cm of y0] {\tiny{$(i\in \Z_{<0})$}};
                \node [main node] (y1) [right = 1.8cm of y0] {};
                \node [main node] (y2) [above left = 0.5cm and 0.5cm of y1] {};
                \node [main node] (y3) [above right = 0.5cm and 0.5cm of y1] {};
                \node [main node] (y4) [below left = 0.5cm and 0.5cm of y1] {};
                \node [main node] (y5) [below left = 0.8cm and 0.8cm of y4] {};
                \node [main node] (y6) [below left = 0.5cm and 0.5cm of y5] {};

                \node (y11) [right = 0.01cm of y1] {\tiny{$0$}};
                \node (y22) [above = 0.01cm of y2] {\tiny{$1$}};
                \node (y33) [above = 0.01cm of y3] {\tiny{$v$}};
                \node (y44) [right = 0.01cm of y4] {\tiny{$-1$}};
                \node (y44) [right = 0.01cm of y5] {\tiny{$(i+1)$}};
                \node (y44) [right = 0.01cm of y6] {\tiny{$i$}};
                \path[draw]
                (y2)--(y1)--(y3)
                (y1)--(y4)
                (y5)--(y6);
                \path[draw,dashed]
                (y4)--(y5);
            \end{tikzpicture}
            &
            \raisebox{1.5em}{
                \begin{tikzpicture}
                    \node (0) {$z_{v}$:};
                    \node [main node] (1) [right = 0.8cm of 0] {};
                    \node [main node] (2) [above left = 0.5cm and 0.5cm of 1] {};
                    \node [main node] (3) [above right = 0.5cm and 0.5cm of 1] {};
                    \node [main node] (4) [below  = 0.5cm  of 1] {};

                    \node (11) [right = 0.01cm of 1] {\tiny{$0$}};
                    \node (22) [above = 0.005cm of 2] {\tiny{$-1$}};
                    \node (33) [above = 0.01cm of 3] {\tiny{$1$}};
                    \node (44) [below = 0.01cm of 4] {\tiny{$v$}};
                    \path[draw]
                    (2)--(1)--(3)
                    (1)--(4);
                \end{tikzpicture}
            }
            & 
            \begin{tikzpicture}
                \node (y0) {$z_i$:};
                \node (b) [below = 0.01cm of y0] {\tiny{$(i\in \Z_{>0})$}};
                \node [main node] (y1) [right = 0.8cm of y0] {};
                \node [main node] (y2) [above left = 0.5cm and 0.5cm of y1] {};
                \node [main node] (y3) [above right = 0.5cm and 0.5cm of y1] {};
                \node [main node] (y4) [below right = 0.5cm and 0.5cm of y1] {};
                \node [main node] (y5) [below right = 0.8cm and 0.8cm of y4] {};
                \node [main node] (y6) [below right = 0.5cm and 0.5cm of y5] {};

                \node (y11) [right = 0.01cm of y1] {\tiny{$0$}};
                \node (y22) [above = 0.005cm of y2] {\tiny{$-1$}};
                \node (y33) [above = 0.01cm of y3] {\tiny{$v$}};
                \node (y44) [right = 0.01cm of y4] {\tiny{$1$}};
                \node (y44) [right = 0.01cm of y5] {\tiny{$(i-1)$}};
                \node (y44) [right = 0.01cm of y6] {\tiny{$i$}};
                \path[draw]
                (y2)--(y1)--(y3)
                (y1)--(y4)
                (y5)--(y6);
                \path[draw,dashed]
                (y4)--(y5);
            \end{tikzpicture}\\\hline
        \end{tabularx}
        \vspace{1em} 
        \caption{The long stubs of $\cals(E_{q,r})$}
        \label{tab:zE}
    \end{table}



We record a few features of the stubs in Theorem~\ref{thm:stubDescription} for future use:

\begin{lemma}
    \label{lemm:stubRemarks}
    Let $(W,S)$ be an arbitrary $\la(2)$-finite Coxeter group from Theorem~\ref{thm:stubDescription} and keep the notation
    of the theorem.  Then the stubs in $\sw$ have the following properties.  
    \begin{enumerate}
        \item The left descent set of each short stub  $x_{tt'}$, medium stub $y_t$ \tul including
             $y'_0,y''_0$ in type $E_{q,r}$\tur and long stub $z_t$ equals $\{t,t'\}, \{t\}$ and
             $\{t\}$, respectively. In particular, a stub $w\in \sw$ has two left descents if and
             only if it is a short stub, and in this case we can recover $w$ from $\call(w)$: if
             $\call(w)=\{t,t'\}$ then $w=x_{tt'}$.
        \item Let $w\in \sw$ be a medium or long stub and let $w=w_p\dots w_1$ be its Cartier--Foata
             factorization.  Then $l(w_1)=2$ and  $l(w_j)=1$ for all $2\le j\le p$, so we may pick a
             reduced word $ \ul{w}= s_1\dots s_{k}(s_{k+1}s_{k+2}) $ of $w$ such that
             $w_1=s_{k+1}s_{k+2}$, $w_2=s_{k}, w_3=s_{k-1}$, etc.
        \item Let $w$ and $\ul w$ be as in \textup{(2)}. Then in the heap $H(\ul w)$, every element
             $i<k+1$ is comparable both to $k+1$ and to $k+2$ via the convex chains of coverings
             $i\prec i+1\prec \dots k\prec k+1$ and $i\prec i+1\prec \dots k\prec k+2$. In
             particular, the element $k$ is both the only element covered by $k+1$ and the only
             element covered by $k+2$ in $H(\ul w)$. 
        \item Let $w$ and $\ul w$ be as in \textup{(2)}. Then we may start with $w$ and successively
             apply left lower star operations with respect to the noncommuting pairs of generators
             $\{s_1,s_2\}, \{s_2,s_3\},\dots, \{s_{k},s_{k+1}\}$ to obtain the sequence $w,
             s_2s_3\dots s_k(s_{k+1}s_{k+2}), s_3\dots s_{k}(s_{k+1}s_{k+2}), \dots,$
             $s_k(s_{k+1}s_{k+2}), w_1$. All elements in this sequence lie in $\sw$ and share the
             same left cell (and hence the same right descents) as $w$. In particular, $w$ is
             in the same left cell and has the same right descents as its first layer $w_1 =
             s_{k+1}s_{k+2}$, which is itself a short stub.
            \end{enumerate}
    \end{lemma}
    \begin{proof} All the claims are readily confirmed by inspection of Theorem~\ref{thm:stubDescription}
          and the pictures in Tables~\ref{tab:xyStubs}--\ref{tab:zE}.     
    \end{proof}

    \begin{remark}
      \label{rmk:a.vs.n}
      The question of whether the $n$-value (as defined in the paragraph before Proposition~\ref{prop:a=n}) and
      $\la$-value of an FC element always agree in a general Coxeter system is open. The results of
      this subsection allow us to enhance Proposition~\ref{prop:a=n}.(1) and give a partial answer to this
      question in the following way: 

\begin{prop} 
\label{prop:a.vs.n}
Let $(W,S)$ be an arbitrary Coxeter system and let $w\in \fc(W)$.  \begin{enumerate} \item We
   have $\la(w)=0$ if and only if $n(w)=0$.  \item We have $\la(w)=1$ if and only if
   $n(w)=1$.
\item If $(W,S)$ is $\la(2)$-finite, then $\la(w)=2$ if and only if $n(w)=2$.  \end{enumerate}
  \end{prop} 
\begin{proof} (1) By Proposition~\ref{prop:facts01}.(1), the equations $\la(w)=0$ and $n(w)=0$ both hold
     exactly when $w$ is the identity element and has an empty reduced word, so they are
     equivalent.

(2) If $\la(w)=1$, then $n(w)\neq 0$ by Part~(1) and $n(w)\le 1$ by Proposition~\ref{prop:a=n}.(1), and
therefore $n(w)=1$.  Conversely, if $n(w)=1$ then the heap $H(w)$ is a chain, so no reduced word
of $w$ contains two consecutive generators that commute. The definition of FC elements then
implies that $w$ has a unique reduced word, so $\la(w)=1$ by Proposition~\ref{prop:facts01}.(2). It follows
that $\la(w)=1$ if and only if $n(w)=1$.

(3) Suppose $(W,S)$ is $\la(2)$-finite. If $\la(w)=2$, then $n(w)\notin\{0,1\}$ by Parts (1) and~(2)
and $n(w)\le 2$ by Proposition~\ref{prop:a=n}.(1), and therefore $n(w)=2$. Now suppose that~$n(w)=2$. We prove that
$\la(w)=2$ by induction on the length $l(w)$ of $w$. We have $l(w)= \abs{H(w)}\ge n(w)=2$, so in the
base case we have $l(w)=2$. Since $n(w)=2$, we must have $w=st$ for two commuting generators $s,t\in
S$ in this case; thus $\la(w)=2$ by Proposition~\ref{prop:Facts}.(2), as desired.  If $l(w)>2$, then $w$
either admits or does not admit a right lower star operation. In the former case, if a right lower
star operation takes $w$ to some element $w'$ then we have $\la(w)=\la(w')=n(w')=n(w)=2$, where the
first equality holds by Corollary~\ref{coro:starA}, the second equality holds by induction, and the third
equality holds because star operations do not change $n$-values of FC elements (see the proof of
Proposition 3.15 in \cite{GreenXu}). In the latter case, the element $w$ is a left stub, so
Corollary~\ref{coro:firstTwoLayers} implies that the first two layers of $w$ satisfy Conditions (1)--(2) from
the corollary. As pointed out in the proof of Theorem~\ref{thm:stubDescription}, this forces $w$ to be an
element in the set $\sw$ given in the theorem, and therefore $w$ indeed has $\la$-value~2.  It
follows by induction that $\la(w)=2$. We have proved that $\la(w)=2$ if and only if~$n(w)=2$.
\end{proof}
    \end{remark}

